\section{Zadanie 15}
\autor{Filip Manijak}
Niech \(X, Y, Z\) będą niezależnymi zmiennymi losowymi o rozkładzie jednostajnym na przedziale \([0, 1]\). Niech \(U = XY\) i \(V = Z^2\).

\begin{enumerate}[label=(\roman*)]
	\item Wyznacz dystrybuantę, gęstość, wartość oczekiwaną i wariancję \(U\).
	\item Wyznacz dystrybuantę, gęstość, wartość oczekiwaną i wariancję \(V\).
	\item Oblicz \(P(U < V)\).
	\item Wyznacz \(E(X | U)\).
\end{enumerate}

\subsection*{Rozwiązanie}
\subsection{U}
Liczymy Dystrybuante:
\[F_U(a) =P(XY\leq a) = \int_0^1f_X(x)\int^{\frac{a}{x}}_0f_Y(y)dy dx = \int_0^af_X(x)*1+ \int_a^1f_X(x) F_Y(\frac{a}{x}) dx =\]
\[ =a+ \int^1_a \frac{a}{x} =a -a \ln a\]
Teraz z tego liczymy funkcję gęstości:
\[f_U(a) = F'_U(a) = (a-a \ln a)' = 1 - \ln a -1 = -\ln a\]
Liczymy też \(E[U]\):
\[E[U] = \int_0^1 f_U(x)xdx = -\int_0^1 x\ln x dx =-\left[ \frac{x^2}{2} \left( \ln x -\frac{1}{2} \right) \right]_0^1  = \frac{1}{4}\]
Zauważmy, że \(\lim_{x\xrightarrow{ }0}x^2 \ln x =0\) bo de l'Hospital. Całkę liczymy całkując przez części. Została wariancja:
\[Var[U] = E[U^2] - E[U]^2 = \int_0^1x^2(-\ln x) dx - \frac{1}{16} =\frac{1}{9} -  \frac{1}{16}\]
Tu też całkę liczymy całkując przez części.
\subsection{V}
Liczymy Dystrybuantę (dla \(a \in (0,1)\)):
\[F_V(a) = P(Z^2 \leq a) = P(-\sqrt{a} \leq Z \leq \sqrt{a}) = \int_0^{\sqrt{a}} 1 dz = \sqrt{a}\]
Teraz z tego liczymy funkcję gęstości:
\[f_V(a) = F'_V(a) = (\sqrt{a})' = \frac{1}{2\sqrt{a}}\]
Liczymy też \(E[V]\):
\[E[V] = \int_0^1 a f_V(a) da = \int_0^1 a \frac{1}{2\sqrt{a}} da = \frac{1}{2} \int_0^1 a^{\frac{1}{2}} da = \frac{1}{2} \left[ \frac{2}{3}a^{\frac{3}{2}} \right]_0^1 = \frac{1}{3}\]
Tutaj całkę liczymy prosto ze wzoru na funkcję potęgową (\(x^n\)). Została wariancja:
\[Var[V] = E[V^2] - E[V]^2 = \int_0^1 a^2 \frac{1}{2\sqrt{a}} da - \frac{1}{9} = \frac{1}{2} \int_0^1 a^{\frac{3}{2}} da - \frac{1}{9} = \frac{1}{5} - \frac{1}{9} = \frac{4}{45}\]
Tu również całkę liczymy jako funkcję potęgową.
\subsection{\(P(U<V)\)}
Zauważmy, że:
\[ P(U< V) = \int_0^1 f_V(x) F_U(x) dx = \int^1_0 \frac{1}{2\sqrt{x}}(x - x\ln x) dx = \frac{1}{2} \left( \int_0^1 \sqrt{x} dx - \int_0^1 \sqrt{x} \ln x dx \right)=\]
\[= \frac{1}{3} - \frac{1}{2}\left[\frac{2}{3}x\sqrt{x}\left(\ln x - \frac{2}{3}\right) \right]^1_0 = \frac{1}{3} - \frac{1}{2} \left( 0 - \frac{4}{9} \right) = \frac{1}{3} + \frac{2}{9} = \frac{5}{9}\]
\subsection{\(E[X | U]\)}
Chcemy obliczyć \(E[X | U](a) = \int_{-\infty}^{\infty} x f_{X|U=a}(x) dx\). W tym celu musimy wyznaczyć gęstość warunkową \(f_{X|U=a}(x)\).
Skorzystamy ze wzoru:
\[f_{X|U=a}(x) = \frac{f_{X,U}(x,a)}{f_U(a)}\]
Mianownik już znamy: \(f_U(a) = -\ln a\). Licznik (gęstość łączną) wyznaczymy, licząc pochodną  z dystrybuanty  \(F_{X,U}(x,a)\).

Dla \(a < x \le 1\):
\[F_{X,U}(x, a) = P(X \le x, XY \le a) = P(X \le a) + P(a < X \le x, Y \le \frac{a}{X}) =\]
\[= \int_0^a 1 dt + \int_a^x \frac{a}{t} dt = a + a(\ln x - \ln a)\]
Liczymy gęstość łączną:
\[f_{X,U}(x, a) = \frac{\partial^2 F_{X,U}}{\partial x \partial a} = \frac{\partial}{\partial a} \left( \frac{\partial}{\partial x} (a + a\ln x - a\ln a) \right) = \frac{\partial}{\partial a} \left( \frac{a}{x} \right) = \frac{1}{x}\]
Zatem gęstość warunkowa dla \(x \in [a, 1]\) wynosi:
\[f_{X|U=a}(x) = \frac{\frac{1}{x}}{-\ln a} = \frac{-1}{x \ln a}\]
Stąd wartość oczekiwana:
\[E[X | U=a] = \int_a^1 x \cdot f_{X|U=a}(x) dx = \int_a^1 x \left( \frac{-1}{x \ln a} \right) dx =\]
\[= \frac{-1}{\ln a} \int_a^1 1 dx = \frac{-1}{\ln a} [x]_a^1 = \frac{-(1-a)}{\ln a} = \frac{a-1}{\ln a}\]
